\ssr{ЗАКЛЮЧЕНИЕ}

Поставленная цель была достигнута: разработан метод построения синтаксических деревьев ограниченного естественного русского языка на основе наиболее употребимых грамматик.

В ходе выполнения выпускной квалификационной работы были выполнены следующие задачи:
\begin{itemize}
	\item проанализирована предметная область автоматизированного построения синтаксических деревьев ограниченного естественного русского языка: выделены и описаны основные формализмы; для реализации выбраны системы составляющих;
	\item проведен сравнительный анализ существующих методов построения синтаксических деревьев: для использования в разрабатываемом методе выбран алгоритм Кока--Янгера--Касами;
	\item сформулированы требования к разрабатываемому методу;
	\item формализована постановка задачи в виде IDEF0-диаграммы;
	\item сформулированы ограничения предметной области;
	\item разработан метод построения синтаксических деревьев ограниченного естественного русского языка;
	\item обоснован выбор программных средств реализации и разработано программное обеспечение реализующее описанный метод;
	\item определена зависимость времени построения синтаксических деревьев от количества словоформ в предложении: наиболее продолжительным этапом является восстановление деревьев составляющих, демонстрирующее экспоненциальный рост временных затрат;
	\item оценена точность и полнота разработанного метода: точность составила 0.68, полнота~---~0.71, что соответствует уровню существующих решений для построения деревьев составляющих на основе готовых деревьев зависимостей.
\end{itemize}

Дальнейшее развитие может быть достигнуто за счет:
\begin{itemize}
	\item более высокой <<детализации>> используемой грамматики, выделения подъединиц текущих синтаксических единиц, ввода новых правил грамматики;
	\item использования более подробного морфологического анализатора, охватывающего все существующие в языке части речи и их морфологические характеристики;
	\item внедрения возможности автоматического разрешения неоднозначности среди всех возвращаемых методом деревьев составляющих;
	\item оптимизации временных затрат этапа восстановления деревьев составляющих.
\end{itemize}



\clearpage